\documentclass{article}
\usepackage{graphicx}
\usepackage{amssymb}
\usepackage{amsmath}
\usepackage{commath}
\usepackage{amsthm}
\usepackage{enumitem}
\usepackage{outlines}
\setlist[enumerate,1]{label=\roman*)}
\setlist[enumerate,2]{label=\alph*)}
\newtheorem{theorem}{Theorem}
\newtheorem*{remark}{Remark}

\usepackage{tikz}
\usetikzlibrary{arrows,shapes,trees}

%%%% PREAMBLE %%%%
\title{Weak solution of the merged mathematical equations of the polluted atmosphere}

%%%% CONTENT %%%%
\begin{document}

\maketitle
\newpage

\begin{abstract}

Abstract

\end{abstract}

{\bf Keywords:} Navier-Stokes equations, shallow domains, pollution evolution equation, hydrostatic approximation.

\newpage

\tableofcontents
\newpage

%%%% INTRODUCTION %%%%
\section{Introduction}

We investigate existence of a weak solution for the merged motion-pollution atmosphere model. At this point the functions we work with are $\mathbf{v}$, and $C$, representing velocity and concentration, respectively (we are not considering other variables such as water vapor quantity or temperature).

%%% *** INLAND DOMAIN *** %%%
\section{Inland domain}

%% Domain Structure
\subsection{Description of the domain and boundary conditions}

A fundamental thing is to precisely decide about and describe the domain and coordinate system we choose to work with.

\begin{itemize}

  \item The set of equations, the choice of the locally acceptable Cartesian coordinate system, and the boundary and initial conditions we are using for the velocity terms are similar to those suggested by the \cite{azerad2001mathematical} article. For describing the effect of pollution we use a continuity equation written in Cartesian coordinates for the concentration. The boundary division of the $\partial \Omega$ domain and the boundary conditions we choose to use are the following.
  
  \begin{figure}[!h]
\centering
\begin{tikzpicture}
\draw (0,0) -- (2,0)node[below] {$\Gamma_G$} -- (4,0) -- (4,2)node[right] {$\Gamma_L$} -- (4,4) -- (2,4)node[above] {$\Gamma_U$} -- (0,4) -- (0,2)node[left] {$\Gamma_L$} -- (0,0);
\draw (2,2)node{$\Omega$};
\draw[green, thick] (0,0) -- (4,0);
\end{tikzpicture}
\end{figure}

  \begin{itemize}
  
      \item At a reasonably high $p = p_0$ pressure level the upper boundary of the local domain is denoted by $\Gamma_U$. On $\Gamma_U$ we use $\bold{u}_H = 0, u_3 = 0, C(x,y,z) = 0$. These zeros are reasonable here since at such a high point it is logical to assume  
    \item On both sides the lateral boundaries of the rectangle shaped $\Omega$ domain are given by $\Gamma_L$. Here we use the same boundary conditions as on the upper level, so we have $\bold{u}_H = 0, u_3 = 0, C(x,y,z) = 0$.
    \begin{remark}
    If we wanted to use $C(x,y,z) = c$ with some $c$ constant value, we could make a change of variables distracting that background flow and arrive back to the same mathematical problem with $C$ being zero.
    \end{remark}
    \item On the $z=0$ ground level we have $u_H = 0, u_3 = 0$ boundary conditions for the air velocity. This is the mathematical way of expressing no-slip boundary conditions for the horizontal wind speed and the impervious nature of the ground with respect to the wind. As for the concentration, the way we can chose ground level boundary conditions is the following. Firstly, since the pollution concentration does not "dig down through" the ground, we can use $\partial_3 C = 0$ on $\Gamma_G$. For the horizontal partial derivatives of the pollution, we can separate downwind and crosswind direction boundary conditions. In the crosswind direction we use $\partial_2 C = 0$ while in the downwind (stronger) direction our assumption is more relaxed, we only assume that $\partial_1 C$ is finite in the boundary norm.
  
  \end{itemize}


\end{itemize}

\subsection{The three main sets of equations}

\subsubsection{The original set of equations}

    \begin{equation}
    \partial_t \bold{v} + (\bold{v} \cdot \nabla)\bold{v} -\Delta_{\nu} \bold{v} + \nabla q + 2 \bold{w} \times \bold{v}= \bold{g} \text{ in } \Omega_{\epsilon} \times (0,T)
    \end{equation}
    
    \begin{equation}
    \nabla \cdot \bold{v} = 0 \text{ in } \Omega_{\epsilon} \times (0,T)
    \end{equation}
    
    \begin{equation}
    \bold{v} = 0 \text{ on } \Gamma^{\epsilon} \times (0,T)
    \end{equation}
    
    \begin{equation}
    \bold{v}(\cdot , t = 0) = \bold{v}_0, C(\cdot , t = 0) = C_0 \text{ in } \Omega_{\epsilon}
    \end{equation}
    
    \begin{equation}
    \partial_t C + \bold{v} \cdot \nabla C = K_x \frac{\partial ^2 C}{\partial x ^ 2} + K_y \frac{\partial ^2 C}{\partial y ^ 2} + K_z \frac{\partial ^2 C}{\partial z ^ 2} + S
    \end{equation}
    
\subsubsection{Rescaling}

    \[x=x_1, y=x_2, z=\epsilon x_3 \]
    \[v_1 = u_1^{\epsilon}, v_2 = u_2^{\epsilon}, v_3 = \epsilon u_3^{\epsilon}, p = p^{\epsilon} \]
    \[\nu_x = \nu_1, \nu_y = \nu_2 , \nu_z = \epsilon ^ 2\nu_3 \]
    \[ \tau ^ \epsilon _ i = \epsilon \theta_i \]
    
    We use an additional scaling equation, which is different in nature from the previous ones in the sense that it is not derived from the shallowness of the domain but rather suggested by the fact that in the downwind direction it is natural to assume that the diffusion is negligible compared to downwind transport,ie.
    
    \begin{equation} \abs{v_1 \partial_x C} \gg \abs{ K_x \frac{\partial ^2 C}{\partial x_1 ^ 2} }\end{equation}
    which leads us to using
    \[ K_x = \epsilon K_1\]


\subsubsection{Anisotropic Navier-Stokes plus pollution}

    \begin{equation}
    \partial_t u^\epsilon _1 + \bold{u}^\epsilon \cdot \nabla u^\epsilon _1 -\Delta_\nu u^\epsilon _1 -\alpha u_2^\epsilon + \epsilon \beta u_3^\epsilon + \partial_1 p^\epsilon = 0 \text{ in } \Omega \times (0,T)
    \end{equation}
    
    \begin{equation}
    \partial_t u_2^\epsilon + \bold{u}^\epsilon \cdot \nabla u_2^\epsilon -\Delta_\nu u_2^\epsilon + \alpha u_1^\epsilon + \partial_2 p^\epsilon = 0 \text{ in } \Omega \times (0,T)
    \end{equation}
    
    \begin{equation}
    \epsilon^2 \{  \partial_t u_3^\epsilon + \bold{u}^\epsilon \cdot \nabla u_3^\epsilon -\Delta_\nu u_3^{\epsilon} \} - \epsilon \beta u_1^\epsilon + \partial_3 p^\epsilon = 0 \text{ in } \Omega \times (0,T)
    \end{equation}
    
     \begin{equation}
    \nabla \cdot \bold{u}^\epsilon= 0 \text{ in } \Omega \times (0,T)
    \end{equation}
    
    \begin{equation}
    \bold{u}^\epsilon = 0 \text{ on } \Gamma \times (0,T)
    \end{equation} 
    
    \begin{equation}
    \bold{u}^\epsilon (\cdot, t = 0) = \bold{u}_0, C^\epsilon (\cdot, t=0) = C^\epsilon_0  \text{ in } \Omega
    \end{equation}
    
    \begin{equation}
    \partial_t C^\epsilon + \bold{u^\epsilon} \cdot \nabla C^\epsilon = \epsilon K_1 \frac{\partial ^2 C^\epsilon}{\partial x_1 ^ 2} + K_2 \frac{\partial ^2 C^\epsilon}{\partial x_2 ^ 2} + K_3 \frac{\partial ^2 C^\epsilon}{\partial x_3 ^ 2} + S
    \end{equation}


\subsubsection{Hydrostatic Navier-Stokes plus pollution}

    \begin{equation}    
    \partial_t u_1 + \bold{u} \cdot \nabla u_1 - \Delta_\nu u_1 -\alpha u_2 + \partial_1 p = 0 \text{ in } \Omega \times (0,T)
    \end{equation}
    
    \begin{equation}    
    \partial_t u_2 + \bold{u} \cdot \nabla u_2 - \Delta_\nu u_2 -\alpha u_1 + \partial_2 p = 0 \text{ in } \Omega \times (0,T)
    \end{equation}
    
    \begin{equation}    
    \partial_3 p = 0 \text{ in } \Omega \times (0,T)
    \end{equation}
    
    \begin{equation}    
    \nabla \cdot \bold{u} = 0 \text{ in } \Omega \times (0,T)
    \end{equation}
    
    \begin{equation}    
    u_1 = u_2 = u_3 n_3 = 0 \text{ on } \Gamma \times (0,T)
    \end{equation}
    
    \begin{equation}    
    \nu_3 \partial_3 u_1 = \theta_1, \nu_3 \partial_3 u_2 = \theta_2 , u_3 = 0 \text{ on } \Gamma_s \times (0,T)
    \end{equation}
    
    \begin{equation}    
    u_i (\cdot, t=0) = u_{0 i}, C (\cdot, t=0) = C_0 \text{ in } \Omega, \text{$i$ = 1,2}
    \end{equation}
    
    \begin{equation}    
        \partial_t C + \bold{u} \cdot \nabla C = K_2 \frac{\partial ^2 C}{\partial x_2 ^ 2} + K_3 \frac{\partial ^2 C}{\partial x_3 ^ 2} + S
    \end{equation}

\subsection{Main theorem}

\begin{itemize}

  
  \item The article of \cite{azerad2001mathematical} proved an existence theorem for weak solution of the equations describing the atmosphere without pollution.
  
  \begin{theorem}
  Let $\bold{u}_0$ $\in$ $L^2 (\Omega)^3$, with $\nabla \cdot \bold{u}_0 = 0, \bold{u}_0 \cdot \bold{n} = 0$ on $\partial \Omega$, and $\theta_1, \theta_2 \in L^2 (0,T; H^(-1/2)(\Gamma_s))$; there exists a weak solution $\bold{u}$ of the hydrostatic Navier-Stokes equations as the aspect ratio $\epsilon$ tends to $0$.
  \end{theorem}
  
  \item The task we want to do is to extend this result in the sense that it becomes valid to the merged set of equations as well.
  
\end{itemize}

\subsection{Weak formulation}

Weak formulation of the merged equations in both the anisotropic and hydrostatic case.

\subsection{Energy inequality and a priori estimates}

\subsection{Passing to the limit}

In the linear terms it is not particularly challenging since we have weak convergence. In the nonlinear term we needed a result that provided a weak convergence of the product of functions for which we only had weak convergence.

%%% *** COASTAL DOMAIN *** %%%
\section{Coastal domain}

%% Domain Structure
\subsection{Domain structure}

Making the model physically more precise, accurate, and adapted for near seaside conditions are the following.
    
\begin{figure}[!h]
\centering
\begin{tikzpicture}
\draw (0,0) -- (2,0)node[below] {$\Gamma_G = \Gamma_E \bigcup \Gamma_I$} -- (4,0) -- (4,2)node[right] {$\Gamma_L$} -- (4,4) -- (2,4)node[above] {$\Gamma_U$} -- (0,4) -- (0,2)node[left] {$\Gamma_L$} -- (0,0);
\draw (2,2)node{$\Omega$};
\draw[green, thick] (0,0) -- (2,0);
\draw[blue,thick,dashed] (2,0) -- (4,0);
\end{tikzpicture}
\end{figure}


On the ground layer we naturally separate the boundary into two different parts, $\Gamma_E$ denoting the part of the boundary above the ground, and $\Gamma_I$ denoting the section that is interacting with the sea. The separation is a sort of expansion or generalization of the ground boundary since while above ground parts it is reasonable to use no-slip conditions, above the ocean it is natural to assume that the wind has an interaction with the waves, there is no such rigidness as in the ground case. For this case we could use $\nu_3 \frac{\partial \bold{u}_H}{\partial x_3} = \theta_H, u_3 = 0$. For the pollution concentration as well, above sea territory it is not natural to use $\partial_3 C = 0$.

%%% *** CONCLUDING REMARKS *** %%%
\section{Concluding remarks, future work}

\begin{itemize}

\item In the paper of \cite{temam2005some}, on p35 it is suggested that the domain we work on has the form of $p_0 < p < p_1$, because for the thin portion of air between the earth and the isobar $p = p_1$, another specific simplified model would be necessary. So a way to make the model more precise is to cut the domain at an isobar a little bit off the physical ground. This would of course radically change the boundary conditions we use for the functions, because it would take away the advantage of having a wall-like, natural and physical boundary condition for our domain to support it. If we are positively above the ground, $\partial_3 C = 0$ would represent a strange reflective layer.

\begin{figure}[!h]
\centering
\begin{tikzpicture}
\draw (0,0) -- (2,0)node[below] {$\Gamma_G = \Gamma_E \bigcup \Gamma_I$} -- (4,0) -- (4,2)node[right] {$\Gamma_L$} -- (4,4) -- (2,4)node[above] {$\Gamma_U$} -- (0,4) -- (0,2)node[left] {$\Gamma_L$} -- (0,0);
\draw (2,2)node{$\Omega$};
\draw[green, thick] (0,-1) -- (2,-1);
\draw[blue,thick,dashed] (2,-1) -- (4,-1);
\end{tikzpicture}
\end{figure}

\item Use more sophisticated ground-level boundary conditions instead of $\partial_3 C = 0$ in the inland case.

\end{itemize}

%%% *** TODO *** %%%
\section{TODO}

1.) Write down the whole process for the version when we touch the ground and there is no ocean

2.) When we touch the ground and there is ocean but no eps in the continuity eq

3.) Look into the physics books

4.) Look into the 2 articles and the idea boundary condition which could be a way for improvement.

\nocite{*}
\bibliographystyle{apalike}
\bibliography{merged_U_C_eqs}

\end{document}
